% =====================================================================
%  Цель и задание лабораторной работы.
% =====================================================================

\textbf{Цель:} получить практические навыки анализа безопасности кода
приложения и состояния его оперативной памяти -- в дополнение к анализу
«чёрного ящика» из предыдущих лабораторных работ -- и на практике
проверить, как ведут себя конфиденциальные и неконфиденциальные данные в
оперативной памяти процесса на всех стадиях их жизненного цикла: при
создании, обновлении и удалении.

\vspace{1em}

\textbf{Задание}:
\begin{enumerate}
  \item Разработать приложение с интерактивным пользовательским
  интерфейсом, позволяющим создать, обновить и удалить данные в/из
  оперативной памяти -- как конфиденциальные, так и неконфиденциальные;
  конфиденциальные данные после ввода должны быть зашифрованы или
  захэшированы.
  \item Выполнить анализ безопасности разработанного приложения с точки
  зрения кода.
  \item Выполнить анализ оперативной памяти и процессорного времени
  приложения.
  \item Создать дамп оперативной памяти приложения после ввода, после
  обновления и после удаления данных; выполнить анализ дампов, уделив
  особое внимание работе с конфиденциальными данными.
  \item Оформить отчёт о проделанной работе.
\end{enumerate}

\textbf{Постановка задачи.} Вместо разработки нового приложения с нуля
существующее приложение \texttt{SecureApp}
(\texttt{Lab2.1-SecureApp}) расширено новой возможностью: разделом
«In-memory data», данные которого умышленно никогда не сохраняются на
диск (не через Entity Framework Core, не через файл, не через
какой-либо кэш) и живут исключительно в памяти процесса backend-сервера
до его перезапуска. Раздел, как и требует задание, поддерживает создание,
обновление и удаление как конфиденциальных данных (шифруемых
непосредственно перед помещением в оперативную память), так и
неконфиденциальных (хранимых в открытом виде) -- через тот же
веб-интерфейс на React, что и постоянные данные лабораторной работы
№2.1, но на отдельных вкладках. Такое расширение позволяет исследовать в
чистом виде именно оперативную память процесса, не смешивая её с
содержимым файла базы данных, которое уже было предметом анализа в
лабораторной работе №2.1.
